\documentclass[fontsize=13pt,paper=b5,leqno,bookmarkpackage=false]{scrartcl}

\PassOptionsToPackage{dvipsnames}{xcolor}
\usepackage[T1]{fontenc}
\usepackage[ngerman,english]{babel}
\usepackage{newpxtext,newpxmath}
\setkomafont{disposition}{\normalcolor\rmfamily\bfseries}
\setkomafont{descriptionlabel}{\normalcolor\rmfamily\bfseries}
\usepackage[
  b5paper,
  left=0.7cm,
  right=0.7cm,
  top=1.0cm,
  bottom=1.0cm,
  footskip=0.55cm
]{geometry}
\linespread{1.1}

\usepackage{microtype}
\usepackage{amsmath,mathtools}
\usepackage{bm}
\usepackage{enumitem}
\usepackage[hang,flushmargin,multiple]{footmisc}
\usepackage{manyfoot}
\usepackage{xcolor}

\definecolor{VNLinkBlue}{HTML}{1F4E79}
\definecolor{VNCiteGreen}{HTML}{2E6B4E}
\definecolor{VNUrlRed}{HTML}{9A3E45}
\definecolor{VNTranslatorPurple}{HTML}{6A3D8E}

\usepackage{hyperref}
\hypersetup{
  colorlinks=true,
  linkcolor=VNLinkBlue,
  citecolor=VNCiteGreen,
  urlcolor=VNUrlRed,
  filecolor=VNTranslatorPurple,
  pdfauthor={J. von Neumann},
  pdftitle={Über adjungierte Funktionaloperatoren / On Adjoint Functional Operators},
  pdfsubject={Faithful English translation of the 1932 paper},
  pdfkeywords={von Neumann, adjoint operators, unbounded operators, polar decomposition},
  bookmarksopen=true,
  bookmarksnumbered=true,
  pdfstartpage=1,
  unicode=true
}
\DeclareNewFootnote{T}
\renewcommand{\thefootnoteT}{T\arabic{footnoteT}}
\makeatletter
\newcommand{\translatornotetarget}{%
  \Hy@raisedlink{%
    \hyper@anchorstart{\Hy@footnote@currentHref}\hyper@anchorend
  }%
}
\makeatother
\newcommand{\translatornote}[1]{%
  \footnoteT{\translatornotetarget
    \textcolor{VNTranslatorPurple}{\textit{Translator's note.}} #1}%
}
\newcommand{\translatornotemark}{\footnotemarkT}
\newcommand{\translatornotetext}[1]{%
  \footnotetextT{\translatornotetarget
    \textcolor{VNTranslatorPurple}{\textit{Translator's note.}} #1}%
}
\newcommand{\fnref}[1]{\hyperref[fn:#1]{note~\ref*{fn:#1}}}
\newcommand{\fnmarkref}[1]{%
  \textsuperscript{\hyperref[fn:#1]{\ref*{fn:#1}}}%
}
\newcommand{\sourcepage}[1]{\par\phantomsection\label{src:#1}}

\setlength{\parindent}{1.35em}
\setlength{\parskip}{0pt}
\setlength{\abovedisplayskip}{7pt plus 2pt minus 2pt}
\setlength{\belowdisplayskip}{7pt plus 2pt minus 2pt}
\setlength{\abovedisplayshortskip}{4pt plus 2pt}
\setlength{\belowdisplayshortskip}{4pt plus 2pt}
\setlist{nosep,leftmargin=2.2em}
\emergencystretch=1.5em
\allowdisplaybreaks[2]

\begin{document}
\selectlanguage{english}
\pdfbookmark[0]{Über adjungierte Funktionaloperatoren / On Adjoint Functional Operators}{bm:title}

% --- Begin inlined translation_pages_01_06.tex ---
% ======================================================================
% Source page 294 (PDF page 1)
% ======================================================================

\phantomsection
\label{article:on-adjoint-functional-operators}
\begingroup
\renewcommand{\thefootnote}{*}
\begin{center}
  {\Large\bfseries Über adjungierte Funktionaloperatoren\footnotemark[0]\par}
  \vspace{0.35em}
  {\Large\bfseries On Adjoint Functional Operators\par}
  \vspace{0.75em}
  {\scshape J. v. Neumann, Princeton\par}
\end{center}
\footnotetext[0]{Received December 9, 1931.}
\endgroup

\section*{Introduction}
\label{sec:introduction}
\addcontentsline{toc}{section}{Introduction}

\phantomsection
\label{intro:1}
\noindent\textbf{1.}
Like several earlier papers,\footnote{\label{fn:1}%
\hypertarget{bib:E}{\href{https://doi.org/10.1007/BF01782338}{%
\emph{Allgemeine Eigenwerttheorie Hermitescher Funktionaloperatoren}}}
[\emph{General Eigenvalue Theory of Hermitian Functional Operators}],
\emph{Math. Annalen} \textbf{102}, no.~1 (1929), pp.~49--131;
\hypertarget{bib:A}{\href{https://doi.org/10.1007/BF01782352}{%
\emph{Zur Algebra der Funktionaloperationen und Theorie
der normalen Funktionaloperatoren}}} [\emph{On the Algebra of Functional Operations and
the Theory of Normal Functional Operators}], \emph{Math. Annalen}
\textbf{102}, no.~3 (1929), pp.~370--427;
\hypertarget{bib:M}{\href{https://doi.org/10.1515/crll.1929.161.208}{%
\emph{Zur Theorie der unbeschränkten Matrizen}}}
[\emph{On the Theory of Unbounded Matrices}], \emph{Journal für Math.}
\textbf{161}, no.~4 (1929), pp.~208--236;
\hypertarget{bib:F}{\href{https://doi.org/10.2307/1968185}{%
\emph{Über Funktionen von Funktionaloperatoren}}}
[\emph{On Functions of Functional Operators}], \emph{Annals of Math.}
\textbf{32}, no.~2 (1931), pp.~191--226. These papers will be cited
repeatedly below, abbreviated respectively as E., A., M., and F.}
this paper is devoted to operators on Hilbert space. The notation and
concept formation follow the author's papers cited in
\hyperref[fn:1]{note~\ref*{fn:1}}; in particular, the terminology to be
used is developed in \hyperlink{bib:E}{E.}, pp.~63--78. Accordingly, an
operator \(A\) is a
function defined on a subset of the Hilbert space \(\mathfrak H\), with
values in \(\mathfrak H\). In general, nothing further is required of an
operator.

The following definitions are essential (cf.\ the work just cited):

\begin{description}
\item[\textbf{I.}]
An operator \(A\) is \emph{defined everywhere} if its domain is all of
\(\mathfrak H\).

\item[\textbf{II.}]
\(A\) is an \emph{extension} of \(B\) if the domain of \(B\) is a subset
of the domain of \(A\), and \(A\) and \(B\) agree on the former; that is,
\(Af=Bf\) there. If their domains are equal, then \(A=B\); otherwise,
\(A\) is a \emph{proper extension} of \(B\).

\item[\textbf{III.}]
\(A\) is \emph{linear} if, together with \(f_1,\ldots,f_n\), all
\(a_1f_1+\cdots+a_nf_n\) belong to its domain,\footnote{\label{fn:2}%
That is, the domain is a linear manifold.} and
\[
  A(a_1f_1+\cdots+a_nf_n)
  =a_1Af_1+\cdots+a_nAf_n .
\]

\item[\textbf{IV.}]
\(A\) is \emph{closed} if, for every sequence \(f_1,f_2,\ldots\) in its
domain such that \(f_n\) converges to some \(f\) and \(Af_n\) converges
to some \(f^*\) as \(n\to\infty\), the element \(f\) also belongs to the
domain,\footnote{\label{fn:3}%
The conclusion that the domain is a closed set would follow from this
only if \(A\) were continuous. Conversely, continuity would follow from
closedness only if \(\mathfrak H\) were compact, which it is not.
(Cf.\ \hyperlink{bib:E}{E.}, p.~70, especially n.~29.)} and \(Af=f^*\).
\end{description}

The operators to be considered will in general not be defined
everywhere. A certain minimum extent of the domain will, however, be
desirable, namely the following:

% ======================================================================
% Source page 295 (PDF page 2)
% ======================================================================

\begin{description}
\item[\textbf{V.}]
\(A\) is a \emph{\(*\)-operator} if the closed linear manifold spanned by
its domain is \(\mathfrak H\),\footnote{\label{fn:4}%
This was required in \hyperlink{bib:E}{E.}, p.~70, Definition~6, for
Hermitian operators, and in \hyperlink{bib:A}{A.}, pp.~405--406,
Definitions~5 and~6, for normal operators.}
that is, if no \(f\ne0\) is orthogonal to its domain.
\end{description}

We now consider the condition
\begin{equation}
  (f,Ag)=(f^*,g),
  \tag*{\(*\)}\label{eq:intro-adjoint-condition}
\end{equation}
where \(f\) and \(f^*\) are regarded as fixed and \(g\) ranges over the
entire domain of \(A\). For a given \(f\), there need not be an \(f^*\)
satisfying \hyperref[eq:intro-adjoint-condition]{\(*\)}; for certain \(f\), however,
there is one, for example \(f=0,\ f^*=0\). The fact that, whenever
\hyperref[eq:intro-adjoint-condition]{\(*\)} is solvable, it has only one solution
\(f^*\) is characteristic of \(*\)-operators.\footnote{\label{fn:5}%
For from one solution \(f^*\) one obtains all the others as \(f^*+h\),
where \(h\) ranges over all elements of \(\mathfrak H\) orthogonal to the
domain of \(A\).} For these, and only for these, the following definition
can therefore be made:

\begin{description}
\item[\textbf{VI.}]
If \(A\) is a \(*\)-operator, then \(A^*\) is the following operator: its
domain consists of those \(f\) for which
\hyperref[eq:intro-adjoint-condition]{\(*\)} is solvable, and in that case
\(A^*f=f^*\).
\end{description}

It is immediately apparent that \(A^*\), which, following the usual
notion that it generalizes,\footnote{\label{fn:6}%
Cf.\ \hyperlink{bib:E}{E.}, pp.~111--112, especially n.~64;
\hyperlink{bib:A}{A.}, p.~372, n.~12, and p.~405, Definition~5.}
is called the \emph{adjoint} of \(A\), is linear
and closed. In general, however, we do not know whether it is defined
anywhere other than at \(0\), still less whether it is a
\(*\)-operator. It is evident that \(A^{**}\), if it can be formed
(that is, if \(A^*\) is a \(*\)-operator), is an extension of \(A\).

These notions permit a simpler description of certain concepts from
\hyperlink{bib:E}{E.}:
a Hermitian operator \(A\) is a \(*\)-operator for which \(A^*\) is an
extension of \(A\); a hypermaximal operator is one for which \(A^*=A\)
(\hyperlink{bib:E}{E.}, p.~70, Definition~6, and p.~72,
Definitions~8 and~9).\footnote{%
\label{fn:7}The equation \(A^*=A\) agrees with the term ``self-adjoint''
used by \hypertarget{bib:Stone1930}{\href{https://doi.org/10.1073/pnas.16.2.172}{%
M.~Stone, \emph{Proc. Nat. Acad.} \textbf{16} (1930), pp.~173--174}},
for hypermaximal operators.}

\phantomsection
\label{intro:2}
\noindent\textbf{2.}
We shall show that a \(*\)-operator \(A\) can be extended to a closed
linear operator \(B\) if and only if \(A^*\) is also a
\(*\)-operator. Among all such \(B\) there is then a smallest one,
that is, one of which every other is an extension; denote it by
\(\widetilde A\). We have
\[
  \widetilde A^{\,*}=A^*,
  \qquad
  \widetilde A=A^{**}.
\]
Another criterion is as follows. Suppose there is a sequence of linear
aggregates
\[
  a_1^{(\nu)}f_1^{(\nu)}+\cdots+
  a_{n_\nu}^{(\nu)}f_{n_\nu}^{(\nu)}
  \qquad
  \bigl(\nu=1,2,\ldots;\ 
  f_1^{(\nu)},\ldots,f_{n_\nu}^{(\nu)}
  \text{ in the domain}\bigr)
\]
such that the aggregates themselves
% ======================================================================
% Source page 296 (PDF page 3)
% ======================================================================
tend to \(0\), while
\[
  a_1^{(\nu)}Af_1^{(\nu)}+\cdots+
  a_{n_\nu}^{(\nu)}Af_{n_\nu}^{(\nu)}
  \longrightarrow f^*
  \qquad(\nu\to\infty).
\]
Then \(f^*=0\). If \(A\) is linear, this has the simpler form: if
\(f_1,f_2,\ldots\) lie in the domain, \(f_n\to0\), and
\(Af_n\to f^*\) as \(n\to\infty\), then \(f^*=0\).\footnote{%
\label{fn:8}Continuity does not follow from this, because the existence
of the limit \(f^*\) was assumed. Cf.\ \hyperref[fn:3]{note~\ref*{fn:3}}.}

These operators are of particular interest:

\begin{description}
\item[\textbf{VII.}]
\(A\) is a \emph{\(**\)-operator} if both \(A\) and \(A^*\) are
\(*\)-operators.
\end{description}

If \(A\) is a \(**\)-operator, then \(\widetilde A\), since it extends
\(A\) and \(\widetilde A^{\,*}=A^*\), is likewise a \(*\)-operator.
Thus \(\widetilde A\) is linear and closed, is a \(**\)-operator, extends
\(A\), and has the same adjoint as \(A\). We therefore obtain a complete
survey of all \(**\)-operators by considering only the closed linear
ones among them. By what was just said, these in turn are identical with
the closed linear \(*\)-operators.

Henceforth we shall therefore always consider closed linear
\(*\)-operators. If \(A\) is such an operator, then, by the foregoing,
\(A^*\) is one as well, and
\[
  A^{**}=A
  \qquad\bigl(\text{since }\widetilde A=A\bigr).
\]

Now consider the operator \(A^*A\), defined as follows: if \(f\) lies in
the domain of \(A\) and \(Af\) lies in the domain of \(A^*\), then
\(A^*Af\) is defined and is equal to \(A^*(Af)\).

From the definition it is clear that \(A^*A0\) is defined and equals
\(0\), and also that \(A^*A\) is linear; its closedness is likewise easy
to prove.\footnote{\label{fn:9}%
From \(f_n\to f\) and \(A^*Af_n\to f^*\) it follows that, for
\(m,n\to\infty\), \(A^*Af_m-A^*Af_n\to0\), and hence
\[
 \begin{aligned}
 &(A^*Af_m-A^*Af_n,f_m-f_n)\\
 &\qquad=(Af_m-Af_n,Af_m-Af_n)
 =\lVert Af_m-Af_n\rVert^2\longrightarrow0 .
 \end{aligned}
\]
Thus \(Af_m\) has a limit, say \(f'\). Since \(A\) is closed, \(Af\) is
defined and equals \(f'\); since \(A^*\) is closed, \(A^*f'\) is defined
and equals \(f^*\) (put \(Af_n=f'_n\)). Consequently \(A^*Af\) is
defined and equals \(f^*\).}
If \(A^*A\) is a \(*\)-operator, it is plainly Hermitian and
definite.\footnote{\label{fn:10}%
\((A^*Af,f)=(Af,Af)=\lVert Af\rVert^2\) is real and nonnegative;
cf.\ \hyperlink{bib:E}{E.}, p.~70, Definition~6, and p.~74,
Definition~10.}
For the moment, however, there is no way to see whether the domain of
\(A^*A\) contains any elements other than \(0\), or even whether it is a
\(*\)-operator. This would be clear only if \(A\) and \(A^*\) were
defined everywhere; but, as we shall show by adapting a familiar
theorem of Toeplitz, that happens only when \(A\) and \(A^*\) are
continuous.

Nevertheless, we shall succeed in proving in full generality that
\(A^*A\) is a \(*\)-operator, indeed a hypermaximal one. Replacing \(A\)
by \(A^*\) shows that the same holds for \(AA^*\).

It is useful to illustrate these relations by some examples. Realize
Hilbert space as the space of all measurable
% ======================================================================
% Source page 297 (PDF page 4)
% ======================================================================
functions \(f(x)\) on the interval \(0\leq x\leq1\) for which
\[
  \int_0^1 |f(x)|^2\,dx<\infty
\]
(cf., for example, \hyperlink{bib:E}{E.}, pp.~108--109). Let
\(\mathfrak A\) be the set of
all continuous differentiable functions \(f(x)\) for which
\[
  \int_0^1 |f'(x)|^2\,dx<\infty,
\]
and which satisfy certain boundary conditions to be specified
immediately; define \(A\) on \(\mathfrak A\) by \(Af=if'\). We allow the
following boundary conditions:
\begin{enumerate}
\item No boundary conditions.
\item \(f(1):f(0)=C\), where \(C\) is a prescribed number.
\item \(f(1)=f(0)=0\).
\end{enumerate}
Correspondingly, we speak of
\[
  \mathfrak A_1,A_1;\qquad
  \mathfrak A_{2,C},A_{2,C};\qquad
  \mathfrak A_3,A_3.
\]
The operators \(A_1,A_{2,C},A_3\) are linear, nonclosed
\(*\)-operators; \(A_1\) extends \(A_{2,C}\), and the latter extends
\(A_3\).\translatornote{The source states the two extension relations in
the opposite direction.  The domains just listed give
\(\mathfrak A_3\subset\mathfrak A_{2,C}\subset\mathfrak A_1\), so the
direction used here follows Definition~II.} No \(A_{2,C}\) is
extended by any other \(A_{2,D}\). From
\phantomsection
\label{eq:intro-integration-by-parts}
\begin{equation*}
 \begin{aligned}
 (Af,g)-(f,Ag)
   &=i\int_0^1
      \left(f'(x)\overline{g(x)}
      +f(x)\overline{g'(x)}\right)\,dx\\
   &=i\bigl(f(x)\overline{g(x)}\bigr)_0^1\\
   &=i\left(f(1)\overline{g(1)}
      -f(0)\overline{g(0)}\right)
 \end{aligned}
\end{equation*}
it follows that \(A_1^*\) is an extension of \(A_3\), likewise
\(A_{2,C}^*\) of \(A_{2,1/\overline C}\), and \(A_3^*\) of \(A_1\).
Thus these three adjoints are \(*\)-operators, and hence
\(A_1,A_{2,C},A_3\) are \(**\)-operators. Consequently,
\phantomsection
\label{eq:intro-differential-examples}
\[
  A_1^*=\widetilde A_1^{\,*},\qquad
  A_{2,C}^*=\widetilde A_{2,C}^{\,*},\qquad
  A_3^*=\widetilde A_3^{\,*}
\]
extend, respectively,
\[
  \widetilde A_3,\qquad
  \widetilde A_{2,1/\overline C},\qquad
  \widetilde A_1.
\]
It is not difficult to show--though we shall not enter into the matter
more closely here--that the respective operators are in fact equal.
It follows that \(\widetilde A_3\) is Hermitian but not hypermaximal;
\(\widetilde A_{2,C}\) is Hermitian only when
\(C=1/\overline C\), that is, \(|C|=1\), but is then also hypermaximal;
\(\widetilde A_1\) is not Hermitian.\translatornote{The source interchanges
the subscripts \(1\) and \(3\) in this sentence.  The adjoint relations
immediately above, together with the boundary conditions, give the
corrected order used here.} On the other hand, the three
operators
\[
 \begin{aligned}
 D_1
   &=\widetilde A_1^{\,*}\widetilde A_1
     =\widetilde A_3\widetilde A_1,\\
 D_{2,C}
   &=\widetilde A_{2,C}^{\,*}\widetilde A_{2,C}
     =\widetilde A_{2,1/\overline C}\widetilde A_{2,C},\\
 D_3
   &=\widetilde A_3^{\,*}\widetilde A_3
     =\widetilde A_1\widetilde A_3
 \end{aligned}
\]
are all hypermaximal. For all three,\translatornote{The source prints
\(Af=-f''\) here; the three operators under discussion are
\(D_1,D_{2,C},D_3\), so \(D_jf=-f''\) is intended.}
\(D_jf=-f''\) \((j=1,2,3)\); only the
domains differ, their boundary conditions being:
\begin{enumerate}
\item For \(D_1\): \(f'(1)=f'(0)=0\).
\item For \(D_{2,C}\):
  \(f(1):f(0)=C,\quad f'(1):f'(0)=1/\overline C\).
\item For \(D_3\): \(f(1)=f(0)=0\).
\end{enumerate}
These are indeed all admissible boundary conditions for the eigenvalue
problem
\[
  -f''=\lambda f,
\]
which is equivalent to hypermaximality (cf.\
\hyperlink{bib:E}{E.}, p.~49, n.~2, and p.~61).

More general or more complicated examples (Sturm--Liouville problems,
partial differential equations, and so forth) exhibit the same
behavior.

\phantomsection
\label{intro:3}
\noindent\textbf{3.}
Since \(A^*A\) and \(AA^*\) are Hermitian, definite, and hypermaximal,
there exists (as will be shown) one and only one Hermitian, definite,
hypermaximal
% ======================================================================
% Source page 298 (PDF page 5)
% ======================================================================
operator \(B\), respectively \(C\), such that
\[
  B^2=A^*A,\qquad C^2=AA^*.
\]
It will be shown that \(B\) has the same domain as \(A\) and that always
\(\lVert Bf\rVert=\lVert Af\rVert\); the same relation holds between
\(C\) and \(A^*\). One can then construct an operator \(W\) which is
essentially length-preserving (cf.\ below) and for which
\phantomsection
\label{eq:intro-polar-relations}
\begin{equation*}
 \begin{gathered}
   A=WB=CW,\qquad
   A^*=BW^*=W^*C,\\
   C=WBW^*,\qquad
   B=W^*CW .
 \end{gathered}
\end{equation*}
As one sees, this is a decomposition of \(A\) and \(A^*\) analogous to
the decomposition of a complex number \(z\) and its conjugate into a
real factor \(\geq0\) (the absolute value) and a factor of absolute
value \(1\). Indeed, in many respects the hypermaximal operators \(B\)
and \(C\) behave like real numbers; definiteness corresponds to
\(\geq0\), and length preservation to absolute value \(1\). The only
striking point is that two different operators, \(B\) and \(C\),
represent the absolute value; they are, however, carried into one
another by the length-preserving mapping \(W\).\footnote{\label{fn:11}%
For completely continuous operators,
\hypertarget{bib:Schmidt1907}{\href{https://doi.org/10.1007/BF01449770}{%
E.~Schmidt}} investigated the eigenvalue problems belonging to \(B\) and
\(C\); cf.\ \emph{Math. Ann.} \textbf{63} (1907), pp.~433--476,
\S\S~14--16.}

For operators \(A,A^*\) defined everywhere (and hence continuous; cf.\
above), this decomposition into \(B,C,W\) is trivial. It is noteworthy,
however, that it persists for every closed linear \(*\)-operator \(A\),
and that \(B\) and \(C\) always turn out to be hypermaximal. Thus, for
example, if \(A\) is Hermitian and definite but not maximal, so that
\(A^*\) is a proper extension of \(A\), and hence
\(A^*A=B^2\) is an extension of \(A^2\), \(B\) is nevertheless
hypermaximal. Thus \(B\ne A\), nor is \(B\) an extension of \(A\), since
it has the same domain!

On the basis of these results it is easy to show that \(B=C\) is
characteristic of the normal operators \(A\), that is, of those
admitting a complex eigenvalue representation
(cf.\ \hyperlink{bib:A}{A.}, p.~406,
Definition~6, and pp.~410--416, Theorems~17--19). Equivalently,
\phantomsection
\label{eq:intro-normality-condition}
\[
  B^2=C^2,
  \qquad\text{or}\qquad
  A^*A=AA^*.
\]
It is well known that this condition characterizes normality in
finite-dimensional space\footnote{\label{fn:12}%
\hypertarget{bib:Frobenius1877}{\href{https://doi.org/10.1515/crelle-1878-18788403}{%
Frobenius, \emph{Journ. f. Math.} \textbf{84} (1877), pp.~51--54}}.}
and for completely continuous operators on Hilbert
space.\footnote{\label{fn:13}%
\hypertarget{bib:HellingerToeplitz}{\href{https://doi.org/10.1007/978-3-663-15917-9}{%
Hellinger and Toeplitz, \emph{Enzykl. d. Math. Wiss.}, II.C.13,
pp.~1562--1563}}.}
The author further showed that, in Hilbert space, it characterizes all
continuous operators (\hyperlink{bib:E}{E.}, pp.~115--116;
\hyperlink{bib:A}{A.}, p.~406; and the work cited
above), but for general \(A\) he obtained there a condition of a
different form. We now recognize \(A^*A=AA^*\) as characteristic in
general; indeed, it is enough that \(A^*A\) be an extension of \(AA^*\),
or conversely, since both operators are hypermaximal.

\phantomsection
\label{intro:4}
\noindent\textbf{4.}
If \(A\) is any linear operator, then on its domain \(\mathfrak A\) we
can introduce a new metric by defining
\[
  \lVert f\rVert_A=\lVert Af\rVert.
\]

% ======================================================================
% Source page 299 (PDF page 6)
% ======================================================================

For \(\mathfrak A\) to be dense in \(\mathfrak H\), \(A\) must be a
\(*\)-operator; for it to be topologically complete in the new metric
(at least relative to \(\mathfrak H\)),\footnote{\label{fn:14}%
That is, \(\lVert f_m-f_n\rVert_A\to0\) as \(m,n\to\infty\) is to imply
the existence in \(\mathfrak A\) of a
\(\lVert\,\cdot\,\rVert_A\)-limit of \(f_1,f_2,\ldots\), insofar as a
\(\lVert\,\cdot\,\rVert\)-limit exists in \(\mathfrak H\). If
\(\lVert Af\rVert\geq\varepsilon\lVert f\rVert\) always holds for some
fixed \(\varepsilon>0\), the latter restriction is superfluous.}
\(A\) must be closed.\footnote{\label{fn:15}%
This consideration of different metrizations in \(\mathfrak H\), as
well as of topological completeness, was suggested by Mr.~K.~Friedrichs
in a lecture that he gave to the Göttingen Mathematical Society in
June 1931.}
Thus we are led back to our class of operators: linear, closed, and
\(*\).

Two operators \(A_1,A_2\) determine the same metric if they have the
same domain and
\(\lVert A_1f\rVert=\lVert A_2f\rVert\) there. We then call them
\emph{metrically equal} and write
\[
  A_1\underset{m}{=}A_2.
\]
If the domain of \(A_1\) is a subset of that of \(A_2\), and
\(\lVert A_1f\rVert=\lVert A_2f\rVert\) on the former, then \(A_2\) is
called a \emph{metric extension} of \(A_1\). If the domains are equal,
then \(A_1\underset{m}{=}A_2\); otherwise, \(A_2\) is a \emph{proper
metric extension} of \(A_1\).

As we have seen, every \(A\) is metrically equal to a Hermitian,
definite, hypermaximal \(B\), and one easily shows that \(B\) is uniquely
determined by this property (cf.\ below).

On the other hand, we shall show that every discontinuous operator \(A\)
has proper metric extensions, even when it is unextendible in the
ordinary sense, that is, maximal. Thus there is no maximality here!

\phantomsection
\label{intro:5}
\noindent\textbf{5.}
Our method will rest essentially on considering the space of all pairs
\(\{f,g\}\), where \(f,g\) are arbitrary elements of \(\mathfrak H\).
If we define
\phantomsection
\label{eq:intro-pair-space}
\begin{equation*}
 \begin{gathered}
   a\{f,g\}=\{af,ag\},\\
   \{f,g\}+\{h,k\}=\{f+h,g+k\},\\
   \left(\!\left(\{f,g\},\{h,k\}\right)\!\right)
      =(f,h)+(g,k),
 \end{gathered}
\end{equation*}
so that
\[
  \lVert\{f,g\}\rVert^2=\lVert f\rVert^2+\lVert g\rVert^2,
\]
this is a Hilbert space, for the characteristic conditions A--E from
\hyperlink{bib:E}{E.}, pp.~64--66, hold for it just as they do for
\(\mathfrak H\). It
stands in approximately the same relation to \(\mathfrak H\) as the
plane does to the line; we shall denote it by \(H\). Every operator \(A\)
on \(\mathfrak H\) can be represented in \(H\) by the set of all pairs
\(\{f,Af\}\). This is the analogue of graphically representing, in the
plane, functions defined on the real line. We shall call this set the
\emph{image} of \(A\), and denote it by \(B_A\).

It is immediately apparent that \(A=B\) means \(B_A=B_B\); that \(A\)
is an extension of \(B\) means that \(B_A\), as a set, contains \(B_B\);
that \(A\) is linear means that \(B_A\) is a linear manifold; and that
\(A\) is closed means that \(B_A\), as a set, is closed.

% --- End inlined translation_pages_01_06.tex ---
% --- Begin inlined translation_pages_07_12.tex ---
% --- Source page 300 (PDF page 7) ---

We shall obtain our results by studying the graphs \(B_A\) in \(H\).
Indeed, the whole theory of operators could be developed in \(H\) much more
systematically than is done here; we shall not, however, pursue this point.

\section*{I. The Operators \(A\), \(A^*\), \(A^*A\).}
\label{sec:operators-aaastara}
\addcontentsline{toc}{section}{\texorpdfstring{I. The Operators \(A\), \(A^*\),
\(A^*A\)}{I. The Operators A, A*, A*A}}

\phantomsection
\label{subsec:I.1}
\noindent\textbf{1.}
As explained in \hyperref[intro:4]{Introduction~4}, we study the Hilbert
space \(H\) and the graphs \(B_A\) in \(H\), in place of the operators \(A\)
in \(\mathfrak H\).  We introduce in \(H\) the further operator
\[
  U\{f,g\}=\{ig,-if\}.
\]
It is Hermitian and unitary, that is,
\(U=U^*=U^{-1}\) (and hence \(U^2=1\)).  Since
\((f,Ag)=(f^*,g)\) can be written as
\[
  \left(\!\left(\{f,f^*\},U\{g,Ag\}\right)\!\right)=0,
\]
the relation \(A^*f=f^*\) says that \(\{f,f^*\}\) is orthogonal to
\(UB_A\).  Thus \(B_{A^*}\) is the set of all elements of \(H\) orthogonal
to \(UB_A\); it is therefore a closed linear manifold, and \(A^*\) is
linear and closed.

The existence of closed linear extensions of \(A\) means that there are
closed linear manifolds \(B_B\) containing \(B_A\).  Now a subset of \(H\)
can be a graph \(B_B\) provided that it never contains two elements
\(\{f,f_1^*\}\), \(\{f,f_2^*\}\) with \(f_1^*\ne f_2^*\); equivalently,
since the subsets in question are linear manifolds, provided that it never
contains an element \(\{0,f^*\}\) with \(f^*\ne0\).  Hence one need only
require that the smallest closed linear manifold containing \(B_A\), namely
the closed linear manifold spanned by \(B_A\), contain no \(\{0,f^*\}\)
with \(f^*\ne0\).  Replacing these notions by their definitions gives:

\par\smallskip
\phantomsection
\label{thm:1}
\noindent\textbf{Theorem 1.}
\textit{The \(*\)-operator \(A\), with domain \(\mathfrak A\), has a closed
linear extension \(B\) if and only if the following condition holds.  If
all}
\[
  f_1^{(n)},\ldots,f_{\nu_n}^{(n)}
  \qquad (n=1,2,\ldots)
\]
\textit{belong to \(\mathfrak A\), and if}
\[
  a_1^{(n)}f_1^{(n)}+\cdots+
  a_{\nu_n}^{(n)}f_{\nu_n}^{(n)}\longrightarrow0,
  \qquad
  a_1^{(n)}Af_1^{(n)}+\cdots+
  a_{\nu_n}^{(n)}Af_{\nu_n}^{(n)}
  \longrightarrow f^*
  \quad(n\to\infty),
\]
\textit{then \(f^*=0\).\footnote{\label{fn:16}The strong convergence used
throughout here could equally well be replaced by weak convergence (see,
for example, \hyperlink{bib:A}{A.}, pp.~378--381), since for linear
manifolds (in this case \(B_A\) and \(B_B\)) strong closedness and weak
closedness are equivalent, as
\hypertarget{bib:Schmidt1908}{\href{https://doi.org/10.1007/BF03029116}{%
E.~Schmidt}} showed in \emph{Rendiconti del Circolo Matematico di Palermo},
\textbf{25} (1908), pp.~57--73.}  Moreover, among all such \(B\) there is a
smallest one, of which all the others are extensions: namely, the operator
for which \(B_B\) is the closed linear manifold spanned by \(B_A\).}

Another criterion is obtained as follows.  If \(A^*\) is also a
\(*\)-operator, then \(A^{**}\) exists; it is linear and closed and is
plainly an extension of \(A\), so that \(\widetilde A\) exists.  Conversely,
if \(\widetilde A\) exists, then \(B_{\widetilde A}\) is the closed linear
manifold spanned by \(B_A\); correspondingly, \(UB_{\widetilde A}\) and
\(UB_A\)
% --- Source page 301 (PDF page 8) ---
have the same orthogonal elements in \(H\), and hence
\(\widetilde A^{\,*}=A^*\).  Furthermore,
\(B_{\widetilde A^{\,*}}=B_{A^*}\) is the orthogonal complement of
\(UB_{\widetilde A}\).  Here two closed linear manifolds are involved; see
\hyperlink{bib:E}{E.}, p.~74.  Since \(U\) is unitary, \(U^2=1\), and
orthogonal complementarity is a symmetric relation, \(B_{\widetilde A}\)
is therefore the orthogonal complement of \(UB_{A^*}\).  In the terminology
of \hyperref[intro:1]{Introduction~1}, this says that Definition~VI there,
applied to \(A^{**}\), yields a unique operator, namely \(\widetilde A\).
Thus \(A^*\) is itself a \(*\)-operator and
\(A^{**}=\widetilde A\).  We have proved:

\par\smallskip
\phantomsection
\label{thm:2}
\noindent\textbf{Theorem 2.}
\textit{The following condition is also necessary and sufficient for the
existence of \(\widetilde A\): \(A^*\) is itself a \(*\)-operator
(\(A\) was already assumed to be one).  In that case}
\[
  \widetilde A^{\,*}=A^*,
  \qquad
  \widetilde A=A^{**}.
\]

It should be noted that no relation whatever between the domains of \(A\)
and \(A^*\) has been required or asserted.

\par\smallskip
\phantomsection
\label{subsec:I.2}
\noindent\textbf{2.}
Let \(A\) now be a closed linear \(*\)-operator.  By
\hyperref[thm:2]{Theorem~2}, \(A^*\) is one as well.  Since \(UB_A\) and
\(B_{A^*}\) are orthogonal complements in \(H\), every element of \(H\)
can be represented, in one and only one way, as the sum of an element of
\(UB_A\) and an element of \(B_{A^*}\) (see, for example,
\hyperlink{bib:E}{E.}, p.~74, Theorem~13).  Choose the element of \(H\) to
be \(\{0,-if\}\).  Then
\[
  \{0,-if\}=\{iAg,-ig\}+\{h,A^*h\}
\]
has a unique solution, with \(g\) in the domain of \(A\) and \(h\) in that
of \(A^*\).  Thus
\[
  iAg+h=0,
  \qquad
  -ig+A^*h=-if.
\]
It follows\translatornote{The
source prints \(Ag=-ih\).  The first displayed equation gives
\(h=-iAg\), which is also the sign used in the following calculation.}
that \(h=-iAg\), equivalently \(Ag=ih\), so \(Ag\) also belongs to the domain of
\(A^*\), and consequently
\[
  -if=-ig+A^*h=-i(g+A^*Ag)=-i(A^*A+1)g,
  \qquad
  (A^*A+1)g=f.
\]
The operator \(A^*A+1\) is therefore defined for sufficiently many \(g\)
that its range is the whole Hilbert space.

Its inverse \(D\), defined by \(Df=g\), is thus linear and everywhere
defined.  Moreover,
\[
\begin{aligned}
  (f,Df)
    &=(A^*Ag+g,g)\\
    &=(A^*Ag,g)+(g,g)\\
    &=(Ag,Ag)+(g,g)
     =\lVert Ag\rVert^2+\lVert g\rVert^2.
\end{aligned}
\]
First, this quantity is real, so \(D\) is Hermitian.  Second, it is
nonnegative and vanishes only when \(\lVert g\rVert^2=0\), that is, when
\(g=0\) and \(f=(A^*A+1)g=0\); hence \(D\) is definite, and \(Df\ne0\)
when \(f\ne0\).  Third, it is at least
\(\lVert g\rVert^2=\lVert Df\rVert^2\), and therefore
\[
  \lVert Df\rVert^2
  \le (f,Df)
  \le \lVert f\rVert\,\lVert Df\rVert,
  \qquad
  \lVert Df\rVert\le\lVert f\rVert.
\]
Thus \(D\) is continuous.  To summarize: \(D\) is a Hermitian, definite,
continuous operator; \(Df\ne0\) whenever \(f\ne0\); its entire spectrum
lies
% --- Source page 302 (PDF page 9) ---
between \(0\) and \(1\),\footnote{\label{fn:17}The last assertion follows
from
\[
  0\le(f,Df)\le\lVert f\rVert^2,
\]
by calculations similar to those in \hyperlink{bib:A}{A.}, p.~418.}
and
\[
  A^*A+1=D^{-1}.
\]
It is now easy to infer the hypermaximality of \(A^*A+1\).  For example,
if \(f\) is orthogonal to every \(Dk\), then \(Df\) is orthogonal to every
\(k\), because \(D\) is Hermitian and everywhere defined.  Hence
\(Df=0\) and \(f=0\).  Thus the orthogonal complement of the closed linear
manifold spanned by the range of \(D\), in \(\mathfrak H\), is \((0)\);
that closed linear manifold is consequently all of \(\mathfrak H\).
The range of \(D\) is the domain of \(A^*A+1\), and hence also the domain
of \(A^*A\), so \(A^*A\) is a \(*\)-operator.  By
\hyperref[intro:2]{Introduction~2}, especially
\fnref{10}, \(A^*A\) is Hermitian and definite, and
therefore so is \(A^*A+1\).  Since \(A^*A+1\) assumes every value in
\(\mathfrak H\), \hyperlink{bib:E}{E.}, p.~102, Theorem~41, applies:
\(A^*A+1\) is hypermaximal, and consequently \(A^*A\) is hypermaximal as
well.  Thus:

\par\smallskip
\phantomsection
\label{thm:3}
\noindent\textbf{Theorem 3.}
\textit{Let \(A\) be a closed linear \(*\)-operator.  Then \(A^*\) is also
a closed linear \(*\)-operator, and \(A^*A\) is Hermitian, definite, and
hypermaximal.}

\par\smallskip
\phantomsection
\label{subsec:I.3}
\noindent\textbf{3.}
The domain \(\mathfrak A_1\) of \(A^*A\) is contained in the domain
\(\mathfrak A\) of \(A\).  Denote by \(A_1\) the restriction of \(A\) to
\(\mathfrak A_1\).  Plainly, \(A_1\) is linear and a \(*\)-operator, but
need not be closed.  Since \(A\) is an extension of \(A_1\), \(A_1^*\) is
an extension of \(A^*\), and is therefore also a \(*\)-operator; hence
\(A_1\) is even a \(**\)-operator.  We shall prove that
\(\widetilde A_1=A\), that is, that \(B_A\) is the closed linear manifold
spanned by \(B_{A_1}\) in \(H\).

Equivalently, the difference between \(B_A\) and the closed linear
manifold spanned by \(B_{A_1}\) is \((0)\).  In other words, the only
element of \(B_A\) orthogonal to that manifold---or, equivalently, to
\(B_{A_1}\)---is \(0\); see \hyperlink{bib:E}{E.}, p.~74,
Definition~11.  An element of \(B_A\) has the form \(\{k,Ak\}\).  If it is
orthogonal to \(B_{A_1}\), then in particular it is orthogonal to the
elements \(\{g,Ag\}\) constructed at the beginning of
\hyperref[subsec:I.2]{I.2}.  Hence \(\{iAk,-ik\}\) is orthogonal to
\(\{iAg,-ig\}\); here \(A^*Ag\) was defined because
\(g\in\mathfrak A_1\).  The element \(\{iAk,-ik\}\), which belongs to
\(UB_A\), is all the more orthogonal to the element \(\{h,A^*h\}\) used
there, since the latter belongs to \(B_{A^*}\).  It is therefore
orthogonal to their sum, \(\{0,-if\}\).  The orthogonality in \(H\) of
\(\{iAk,-ik\}\) and \(\{0,-if\}\) means that \(k\) and \(f\) are
orthogonal in \(\mathfrak H\).  Since this holds for every \(f\), we have
\(k=0\) and \(\{k,Ak\}=\{0,0\}\).  This proves the assertion:

\begin{samepage}
\par\smallskip
\phantomsection
\label{thm:4}
\noindent\textbf{Theorem 4.}
\textit{Let the domain of \(A\) be \(\mathfrak A\), and let the domain of
\(A^*A\) be \(\mathfrak A_1\subset\mathfrak A\).  Denote by \(A_1\) the
restriction of \(A\) to \(\mathfrak A_1\).  Then \(A_1\) is a linear
\(**\)-operator, though it need not be closed.  Nevertheless,}
\[
  \widetilde A_1=A.
\]
\end{samepage}

\par\smallskip
\phantomsection
\label{subsec:I.4}
\noindent\textbf{4.}
If \(B\) is a Hermitian hypermaximal operator, it has an associated
resolution of the identity \(E(\lambda)\), by \hyperlink{bib:E}{E.},
p.~91, Definition~17, and p.~92, Theorem~36.  If \(E(\lambda)\) is
constant, and hence equal to \(0\), for \(\lambda<0\), the formulas there
give \((Bf,f)\ge0\) for every \(f\); thus \(B\) is definite.  If, on the
other hand, \(E(\lambda)\ne0\) for some \(\lambda<0\), then for some such
\(\lambda=\lambda_0<0\) there is an \(f\ne0\) with
\(E(\lambda_0)f=f\),
% --- Source page 303 (PDF page 10) ---
and therefore \(E(\lambda)f=f\) for \(\lambda\ge\lambda_0\), since
\(E(\lambda_0)\) is contained in \(E(\lambda)\); compare the discussion
of projection operators in \hyperlink{bib:E}{E.}  From the same formulas
one then obtains
\[
  (Bf,f)\le\lambda_0\lVert f\rVert^2<0,
\]
so \(B\) is not definite.  Thus \(B\) is definite if and only if
\(E(\lambda)\) is constant, and hence \(0\), for \(\lambda<0\).

Now consider the family\translatornotemark
\[
  F(\mu)=
  \begin{cases}
    0, & \mu<0,\\
    E(\sqrt{\mu}), & \mu\ge0.
  \end{cases}
\]
\translatornotetext{The source prints \(E(\sqrt{\lambda})\) here; the subsequent
calculation and the definition as a function of \(\mu\) require
\(E(\sqrt{\mu})\), used above.}
It too is plainly a resolution of the identity, and its associated
Hermitian hypermaximal operator \(B'\) is definite.  If \(B'f\) is
defined, then, by definition,
\[
  \int \mu^2\,d\lVert F(\mu)f\rVert^2
  =
  \int \mu^2\,d\lVert E(\sqrt{\mu})f\rVert^2
  =
  \int \lambda^4\,d\lVert E(\lambda)f\rVert^2
\]
is finite.  Hence
\(\int\lambda^2\,d\lVert E(\lambda)f\rVert^2\) is finite as well, for
\(\int d\lVert E(\lambda)f\rVert^2\) is also finite and
\[
  \lambda^2\le \tfrac12+\tfrac12\lambda^4.
\]
Thus \(Bf\) is defined.  The same calculation carried out in
\hyperlink{bib:F}{F.}, p.~206, to prove formula~e) there shows that
\[
  \int\lambda^2\,d\lVert E(\lambda)Bf\rVert^2
  =
  \int\lambda^4\,d\lVert E(\lambda)f\rVert^2.
\]
This is finite, so \(B(Bf)\), and hence \(B^2f\), is defined.  One more
calculation of the same kind gives
\[
\begin{aligned}
  (B^2f,g)
    &=(B(Bf),g)
      =\int\lambda\,d(E(\lambda)Bf,g)
      =\int\lambda^2\,d(E(\lambda)f,g)\\
    &=\int\mu\,d(E(\sqrt{\mu})f,g)
      =\int\mu\,d(F(\mu)f,g)
      =(B'f,g).
\end{aligned}
\]
Thus \(B^2f=B'f\).  Hence \(B^2\) is an extension of the hypermaximal
operator \(B'\), and therefore \(B^2=B'\).

Rearranging the statement slightly, we may formulate it as follows:

\par\smallskip
\phantomsection
\label{thm:5}
\noindent\textbf{Theorem 5.}
\textit{Let \(B\) and \(B'\) be Hermitian, hypermaximal, definite
operators, and let \(E(\lambda)\) and \(F(\lambda)\), respectively, be
their resolutions of the identity; thus
\(E(\lambda)=F(\lambda)=0\) for \(\lambda<0\).  Then}
\[
  B^2=B'
\]
\textit{is equivalent to}
\[
  E(\lambda)=F(\lambda^2)
  \qquad\textit{for every }\lambda\ge0.
\]

Consequently, if \(B'\) is given and \(B\) is not, the equation
\(B^2=B'\) has exactly one solution \(B\) of this kind; it will be denoted
by \(\sqrt{B'}\).

We now apply \hyperref[thm:5]{Theorem~5} to \(B'=A^*A\), and set
\[
  B=\sqrt{A^*A}.
\]
The domain of \(B^2=A^*A\) was denoted by \(\mathfrak A_1\); let \(A_1\)
be the restriction of \(A\) to this domain,
% --- Source page 304 (PDF page 11) ---
and let \(B_1\) be the restriction of \(B\) to it.  The domain of \(B\)
contains that of \(B^2\), namely \(\mathfrak A_1\).  By
\hyperref[thm:4]{Theorem~4},\translatornote{The source prints
``Theorem~3'' here, but the assertion \(\widetilde A_1=A\) is
\hyperref[thm:4]{Theorem~4}, which is plainly intended.}
\(\widetilde A_1=A\); replacing \(A\) there
by \(B\), and using \(B^*=B\) by hypermaximality, gives
\(\widetilde B_1=B\).

The operators \(A_1\) and \(B_1\) have the same domain, and on that domain
\[
\begin{aligned}
  \lVert A_1f\rVert^2
    &=(A_1f,A_1f)=(Af,Af)=(A^*Af,f),\\
  \lVert B_1f\rVert^2
    &=(B_1f,B_1f)=(Bf,Bf)=(B^2f,f).
\end{aligned}
\]
Thus \(\lVert A_1f\rVert=\lVert B_1f\rVert\).  Hence \(A_1\) and \(B_1\)
are metrically equal (see \hyperref[intro:4]{Introduction~4}), and it
follows that their closures \(\widetilde A_1\) and
\(\widetilde B_1\)\footnote{\label{fn:18}Let \(C\) be a linear
\(**\)-operator (by \hyperref[thm:4]{Theorem~4}, \(A_1\) and \(B_1\) are
such operators).  Then \(B_C\) is a linear manifold and
\(B_{\widetilde C}\) is its closure.  In other words, \(\widetilde C\) is
defined as follows: if \(f_1,f_2,\ldots\) belong to the domain of \(C\),
if \(f_n\to f\) as \(n\to\infty\), and if the \(Cf_n\) also have a limit,
that is,
\[
  \lVert C(f_m-f_n)\rVert\longrightarrow0
  \qquad(m,n\to\infty),
\]
then \(f\) belongs to the domain of \(\widetilde C\), and
\(\widetilde Cf\) is the limit of the \(Cf_n\).  This formulation makes
it clear that, when \(C\) is replaced by a metrically equal operator,
\(\widetilde C\) is likewise replaced by a metrically equal operator.}
are metrically equal as well; that is, \(A\) and \(B\) are metrically
equal.  Substitution into the definition of metric equality yields:

\par\smallskip
\phantomsection
\label{thm:6}
\noindent\textbf{Theorem 6.}
\textit{Let \(A\) be as in \hyperref[thm:2]{Theorems~2} and
\hyperref[thm:3]{3}, and let}
\[
  B=\sqrt{A^*A}
  \qquad\textit{(see \hyperref[thm:5]{Theorem~5})}.
\]
\textit{Thus \(B\) is Hermitian, definite, and hypermaximal.  Then \(A\)
and \(B\) are metrically equal: they have the same domain
\(\mathfrak A\), and throughout that domain}
\[
  \lVert Af\rVert=\lVert Bf\rVert.
\]

\par\smallskip
\phantomsection
\label{subsec:I.5}
\noindent\textbf{5.}
Denote the range of \(A\) by \(\mathfrak W\), and denote the closures of
the linear manifolds \(\mathfrak A\) and \(\mathfrak W\), respectively,
by
\[
  \overline{\mathfrak A}=\mathfrak H,
  \qquad
  \overline{\mathfrak W}.
\]
Further, let \(\mathfrak N\) be the set of all \(f\) such that \(Af=0\).
It is contained in \(\mathfrak A\), and hence in
\(\overline{\mathfrak A}\), and, since \(A\) is linear and closed, it is
a closed linear manifold (all these sets being in \(\mathfrak H\)).
For \(A^*\), denote the corresponding sets by
\[
  \mathfrak A',\quad
  \mathfrak W',\quad
  \overline{\mathfrak A}'=\mathfrak H,\quad
  \overline{\mathfrak W}',\quad
  \mathfrak N'.
\]

For \(f\) to be orthogonal to \(\overline{\mathfrak W}\) is the same as
for \(f\) to be orthogonal to \(\mathfrak W\), that is, to every \(Ag\).
Thus \((f,Ag)=0\) for every \(g\), which says precisely that \(A^*f\) is
defined and equal to \(0\), or that \(f\in\mathfrak N'\).  Therefore
\(\mathfrak N'=\mathfrak H-\overline{\mathfrak W}\).  Similarly,
\(\mathfrak N=\mathfrak H-\overline{\mathfrak W}'\).  In summary,
\[
  \mathfrak N\subset\mathfrak A,
  \qquad
  \mathfrak N=\mathfrak H-\overline{\mathfrak W}',
  \qquad
  \mathfrak N'\subset\mathfrak A',
  \qquad
  \mathfrak N'=\mathfrak H-\overline{\mathfrak W}.
\]

By \hyperref[thm:6]{Theorem~6},\translatornote{The source prints
``Theorem~4'' here; the equality of the two null sets follows from the
metric equality asserted in \hyperref[thm:6]{Theorem~6}.} the set on which
\(Af=0\) is identical
with the set on which \(Bf=0\); hence \(\mathfrak N\) is also the
corresponding set for \(B\).  Let \(\mathfrak B\) be the range of \(B\),
a linear manifold, and let \(\overline{\mathfrak B}\) be its closure, a
closed linear manifold.  The same argument gives
\[
  \mathfrak N=\mathfrak H-\overline{\mathfrak B},
\]
because \(B^*=B\) by hypermaximality.  Consequently,
\[
  \overline{\mathfrak B}=\overline{\mathfrak W}'.
\]

Replace \(A\) by \(A^*\), and hence \(A^*\) by \(A^{**}=A\).
By \hyperref[thm:5]{Theorem~5} and \hyperref[thm:6]{Theorem~6}, this
produces the Hermitian, definite, hypermaximal operators \(AA^*\) and
\[
  C=\sqrt{AA^*}.
\]
By \hyperref[thm:6]{Theorem~6}, \(A^*\) and \(C\) are metrically equal:
they have the same domain \(\mathfrak A'\), and on it
\[
  \lVert A^*f\rVert=\lVert Cf\rVert.
\]

% --- Source page 305 (PDF page 12) ---

For \(C\), let \(\mathfrak B'\) be its range and
\(\overline{\mathfrak B}'\) its closure.  Then
\[
  \overline{\mathfrak B}'=\overline{\mathfrak W}.
\]
By \hyperref[thm:6]{Theorem~6},\translatornote{The source again prints
``Theorem~4'' here; the isometry of the correspondence is the conclusion
of \hyperref[thm:6]{Theorem~6}.} the correspondence
\[
  Bf\Longleftrightarrow Af
\]
maps \(\mathfrak B\) one-to-one, linearly, and isometrically onto
\(\mathfrak W\).  It can therefore be extended, again one-to-one,
linearly, and isometrically, to the closures
\[
  \overline{\mathfrak B}=\overline{\mathfrak W}'
  \quad\hbox{and}\quad
  \overline{\mathfrak W},
\]
that is, from
\(\mathfrak H-\mathfrak N\) onto
\(\mathfrak H-\mathfrak N'\).  Denote the operator of this isometric
mapping by \(V\).  Its domain is
\(\mathfrak H-\mathfrak N\), and its range is
\(\mathfrak H-\mathfrak N'\); compare
\hyperlink{bib:E}{E.}, p.~70, Definition~6, and p.~71, Theorem~10.

Similarly,
\[
  Cf\Longleftrightarrow A^*f
\]
is a mapping of \(\mathfrak B'\) onto \(\mathfrak W'\), and it can be
extended to \(\overline{\mathfrak B}'=\overline{\mathfrak W}\) and
\(\overline{\mathfrak W}'\), that is, from
\(\mathfrak H-\mathfrak N'\) onto
\(\mathfrak H-\mathfrak N\).  Denote its isometric operator by
\(V'\); its domain is \(\mathfrak H-\mathfrak N'\), and its range
is \(\mathfrak H-\mathfrak N\).  By the definitions of \(V\) and
\(V'\),
\[
  A=VB,
  \qquad
  A^*=V'C.
\]

The way in which \(V\), respectively \(V'\), is defined outside the range
of \(B\), respectively \(C\), that is, outside
\(\mathfrak H-\mathfrak N\), respectively
\(\mathfrak H-\mathfrak N'\), is immaterial.  We may therefore
extend them to everywhere-defined linear operators \(W\), respectively
\(W'\), which vanish on \(\mathfrak N\), respectively \(\mathfrak N'\).
Then again
\[
  A=WB,
  \qquad
  A^*=W'C.
\]

Let \(E\), respectively \(E'\), denote the projection operators onto
\(\mathfrak H-\mathfrak N\), respectively
\(\mathfrak H-\mathfrak N'\).  Then
\[
  W=VE,
  \qquad
  W'=V'E',
\]
and therefore immediately
\[
  \lVert Wf\rVert=\lVert Ef\rVert,
  \qquad
  \lVert W'f\rVert=\lVert E'f\rVert.
\]
This proves that \(W\) and \(W'\) are continuous.  It also gives
\begin{equation}
  \tag*{\(\circledcirc\)}
  \label{eq:WstarW}
  W^*W=E,
  \qquad
  W'^{*}W'=E'.\footnotemark
\end{equation}
\footnotetext{\label{fn:19}Indeed,
  \[
    \lVert Wf\rVert^2=(Wf,Wf)=(W^*Wf,f),
    \qquad
    \lVert Ef\rVert^2=(Ef,f),
  \]
  and hence \(W^*W=E\); similarly \(W'^{*}W'=E'\).}

Since \(W\) maps\translatornote{The source prints \(W'\) here.  The mapping
described is \(W\), as required by \(W^*W=E\) and by the domain and range
stated in the same sentence.}
\(\mathfrak H-\mathfrak N\) onto
\(\mathfrak H-\mathfrak N'\), while \(W^*W=E\) is the identity
mapping of \(\mathfrak H-\mathfrak N\) onto itself, \(W^*\) effects
the inverse, likewise isometric, mapping from
\(\mathfrak H-\mathfrak N'\) onto
\(\mathfrak H-\mathfrak N\).  Hence
\(\lVert W^*f\rVert=\lVert f\rVert\) on
\(\mathfrak H-\mathfrak N'\).  On \(\mathfrak N'\), since it is
orthogonal to the range \(\mathfrak H-\mathfrak N'\) of \(W\), one
has \(W^*f=0\).  Thus, in general, \(W^*E'f=W^*f\), and therefore
\[
  \lVert W^*f\rVert=\lVert E'f\rVert.
\]
It follows (compare \fnref{19}) that
\begin{equation}
  \tag*{\(\circledcirc\)}
  \label{eq:WWstar}
  WW^*=E',
\end{equation}

% --- End of source page 305; the argument continues on source page 306. ---

% --- End inlined translation_pages_07_12.tex ---
% --- Begin inlined root_pages_13_14_backup.tex ---
% --- Source page 306 (PDF page 13) ---

and likewise
\begin{equation}
  \tag*{\(\circledcirc\)}
  \label{eq:WpWpstar}
  W'W'^{*}=E.
\end{equation}

Left multiplication of \(A=WB\) by \(W^*\) gives, since \(EB=B\)
(the values of \(B\) lie in \(\mathfrak H-\mathfrak N\)),
\begin{equation}
  \tag*{\(*\)}
  \label{eq:star-B}
  B=W^*A.
\end{equation}
Likewise, left multiplication of \(A^*=W'C\) by \(W'^{*}\) gives
\begin{equation}
  \tag*{\(*\)}
  \label{eq:star-C}
  C=W'^{*}A^*.
\end{equation}
Taking adjoints in these four relations yields, after a simple
discussion,
\begin{equation}
  \tag*{\(\#\)}
  \label{eq:hash-adjoints}
  A^*=BW^*,
  \qquad
  A=CW'^{*},
\end{equation}
\begin{equation}
  \tag*{\(\#\)}
  \label{eq:hash-converses}
  B=A^*W,
  \qquad
  C=AW'.
\end{equation}

We now form the operator \(WBW^*\).\footnote{\label{fn:20}%
Operator products are consistently understood as follows: \(RSf\) is
defined when \(Sf\) and \(R(Sf)\) are defined, and it is equal to
\(R(Sf)\); correspondingly for \(QRS,PQRS,\ldots\).  This multiplication
is plainly associative.}
Since \(W\) is defined everywhere, this product is defined whenever
\(W^*f\) belongs to \(\mathfrak A\).  Since
\(\mathfrak N\subset\mathfrak A\), the projection onto \(\mathfrak N\)
of every element of \(\mathfrak A\) belongs to \(\mathfrak A\), and so
does its projection onto \(\mathfrak H-\mathfrak N\).  Thus the part of
\(\mathfrak A\) lying in \(\mathfrak H-\mathfrak N\) is everywhere
dense there; equivalently, the part of \(\mathfrak A\) lying in the
range of \(W^*\), which is \(\mathfrak H-\mathfrak N\), is everywhere
dense there.  Since \(W^*\) maps
\(\mathfrak H-\mathfrak N'\) one-to-one and isometrically onto
\(\mathfrak H-\mathfrak N\), those \(f\) in
\(\mathfrak H-\mathfrak N'\) for which \(W^*f\in\mathfrak A\) are
everywhere dense in that subspace.  For \(f\in\mathfrak N'\) one has
\(W^*f=0\), which also belongs to \(\mathfrak A\).  Hence the \(f\) just
described are everywhere dense in \(\mathfrak H\), and \(WBW^*\) is a
\(*\)-operator.

It is Hermitian and definite just as \(B\) is.  It is also hypermaximal.
Indeed, suppose that, for every \(g\) for which the expression is
defined,
\[
  (f,WBW^*g)=(f^*,g),
\]
or equivalently
\[
  (W^*f,BW^*g)=(f^*,g).
\]
Put \(g=Wh\), where \(h\in\mathfrak H-\mathfrak N\) and \(Bh\) is
defined.  Then \(W^*g=W^*Wh=h\), and therefore
\[
  (W^*f,Bh)=(f^*,Wh)=(W^*f^*,h).
\]
For every \(h\in\mathfrak N\), \(Bh\) is defined and equals \(0\); the
right-hand side also vanishes because \(W^*f^*\) belongs to
\(\mathfrak H-\mathfrak N\).  Consequently,
\[
  (W^*f,Bh)=(W^*f^*,h)
\]
holds for every \(h\) in the domain of \(B\).

% --- Source page 307 (PDF page 14) ---

Since \(B\) is hypermaximal, \(BW^*f\) is defined and equals \(W^*f^*\).
Thus \(WBW^*f=WW^*f^*=f^*\), provided
\(f^*\in\mathfrak H-\mathfrak N'\).\translatornote{The source prints
\(f^*\in\mathfrak N'\); the argument immediately following shows that
\(f^*\) is orthogonal to \(\mathfrak N'\), hence belongs to
\(\mathfrak H-\mathfrak N'\).} This last condition follows from
the original equation: for every \(g\in\mathfrak N'\), \(W^*g=0\), so
its left-hand side is \(0\), and hence \((f^*,g)=0\).

Now, using the
\hyperref[eq:star-B]{\(*\)},
\hyperref[eq:WstarW]{\(\circledcirc\)}, and
\hyperref[eq:hash-adjoints]{\(\#\)} formulas\fnmarkref{20}, one obtains
\phantomsection
\label{eq:WBWstar-square}
\[
\begin{aligned}
  (WBW^*)^2
    &=WBW^*\cdot WBW^*
      =WB\cdot W^*W\cdot BW^*\\
    &=WB\cdot E\cdot BW^*
      =WB\cdot EB\cdot W^*\\
    &=WB\cdot BW^*
      =A\cdot A^*
      =C^2.
\end{aligned}
\]
Therefore, by \hyperref[thm:5]{Theorem~5},\translatornote{The source
prints ``Theorem~4'' both here and in the parenthetical reference in
\hyperref[thm:7]{Theorem~7}.  Here the intended result is plainly
\hyperref[thm:5]{Theorem~5}, the uniqueness of the definite hypermaximal
square root; in Theorem~7 the intended result is
\hyperref[thm:6]{Theorem~6}, applied to \(A\) and \(A^*\).}
\[
  WBW^*=C.
\]
By \hyperref[eq:hash-adjoints]{\(\#\)}, this says \(WA^*=C\), while
\hyperref[eq:star-C]{\(*\)} says \(W'^{*}A^*=C\).  Thus
\(Wf=W'^{*}f\) throughout the range \(\mathfrak W'\) of \(A^*\), and
hence also throughout
\(\overline{\mathfrak W}'=\mathfrak H-\mathfrak N\).  On
\(\mathfrak N\) the equality holds as well, since both sides vanish.
Consequently \(W=W'^{*}\); taking adjoints gives \(W^*=W'\).

We may summarize the results as follows.

\par\smallskip
\phantomsection
\label{thm:7}
\noindent\textbf{Theorem 7.}
\textit{Let \(A\) be as in \hyperref[thm:2]{Theorems~2} and
\hyperref[thm:3]{3}, and set}
\[
  B=\sqrt{A^*A},
  \qquad
  C=\sqrt{AA^*}.
\]
\textit{(Compare \hyperref[thm:6]{Theorem~6}, applied to \(A\) and to
\(A^*\).)}
\textit{Thus \(B\) and \(C\) are Hermitian, definite, and hypermaximal.
Let the domains of \(A,A^*,B,C\) be, respectively, the linear manifolds}
\[
  \mathfrak A,\quad \mathfrak A',\quad
  \mathfrak A,\quad \mathfrak A',
\]
\textit{and let their ranges be, respectively,}
\[
  \mathfrak W,\quad \mathfrak W',\quad
  \mathfrak B,\quad \mathfrak B',
\]
\textit{with an overline denoting closure.  Let the respective sets
(closed linear manifolds) on which \(Af=0,A^*f=0,Bf=0,Cf=0\) be}
\[
  \mathfrak N,\quad \mathfrak N',\quad
  \mathfrak N,\quad \mathfrak N'.
\]
\textit{Then}
\[
\begin{gathered}
  \mathfrak N\subset\mathfrak A\subset
  \overline{\mathfrak A}=\mathfrak H,
  \qquad
  \mathfrak N'\subset\mathfrak A'\subset
  \overline{\mathfrak A}'=\mathfrak H,\\
  \overline{\mathfrak W}'=\overline{\mathfrak B}
    =\mathfrak H-\mathfrak N,
  \qquad
  \overline{\mathfrak W}=\overline{\mathfrak B}'
    =\mathfrak H-\mathfrak N'.
\end{gathered}
\]
\textit{Let \(E,E'\) be the projection operators onto
\(\mathfrak H-\mathfrak N\) and
\(\mathfrak H-\mathfrak N'\), respectively.  There exists an
everywhere-defined continuous operator \(W\) which maps
\(\mathfrak B\) one-to-one, linearly, and isometrically onto
\(\mathfrak W\), and therefore maps
\(\overline{\mathfrak B}=\mathfrak H-\mathfrak N\) onto
\(\overline{\mathfrak W}=\mathfrak H-\mathfrak N'\).
Likewise \(W^*\) maps \(\mathfrak B'\) onto \(\mathfrak W'\), and
therefore maps
\(\overline{\mathfrak B}'=\mathfrak H-\mathfrak N'\) onto
\(\overline{\mathfrak W}'=\mathfrak H-\mathfrak N\); on these latter
subspaces \(W^*\) and \(W\) are inverse to one another.  The operator
\(W\) vanishes on \(\mathfrak N\), and \(W^*\) on
\(\mathfrak N'\).  Moreover,}
\[
  W^*W=E,
  \qquad
  WW^*=E'.
\]
\textit{The following relations hold\fnmarkref{20}:}
\phantomsection
\label{eq:theorem7-relations}
\[
\begin{gathered}
  A=WB=CW,
  \qquad
  A^*=BW^*=W^*C,\\
  B=W^*A=A^*W,
  \qquad
  C=AW^*=WA^*,
\end{gathered}
\]
\textit{and consequently}
\[
  C=WBW^*,
  \qquad
  B=W^*CW.
\]

% --- End inlined root_pages_13_14_backup.tex ---
% --- Begin inlined translation_pages_15_17_backup.tex ---
% --- Source page 308 (PDF page 15) ---

It should be noted that \(B\) vanishes on \(\mathfrak N\), and \(C\) on
\(\mathfrak N'\), so that they are of interest only on
\(\mathfrak H-\mathfrak N\) and \(\mathfrak H-\mathfrak N'\),
respectively.  On these subspaces, however, they pass into one another
under the mutually inverse isometric mappings between
\(\mathfrak H-\mathfrak N\) and \(\mathfrak H-\mathfrak N'\) effected by
\(W\) and \(W^*\).  This is the generalization of E.~Schmidt's theorem
mentioned in \hyperref[intro:3]{Introduction~3} and \fnref{11},
respectively.

If \(0\) is a (point-)eigenvalue of neither \(A\) nor \(A^*\), that is, if
\(Af=0\) and \(A^*f=0\) each imply \(f=0\), then
\[
  \mathfrak N=\mathfrak N'=(0),
  \qquad
  E=E'=1,
\]
and hence \(W^*W=WW^*=1\); that is, \(W\) is unitary and
\(W^*=W^{-1}\).  In this case the formulas of
\hyperref[thm:7]{Theorem~7} are still more transparent.

\section*{II. Applications.}
\label{sec:applications}
\addcontentsline{toc}{section}{II. Applications}

\phantomsection
\label{subsec:II.1}
\noindent\textbf{1.}
If \(A^*A=AA^*\), that is, if \(B=C\), then
\hyperref[thm:7]{Theorem~7} gives
\[
  A=WB=BW,
  \qquad
  A^*=BW^*=W^*B,
\]
\begin{samepage}
and also\translatornotemark
\[
  \mathfrak N=\mathfrak N',
  \qquad
  E=E',
  \qquad
  W^*W=WW^*=E.
\]
\translatornotetext{The source prints \(E=F\); in the notation of
\hyperref[thm:7]{Theorem~7}, \(E'\) is intended.}
\end{samepage}
Thus all the operators \(B,W,W^*\) commute with one another.  From this
the normality of \(A\) (compare \hyperref[intro:3]{Introduction~3}) may be
deduced in various ways.  For example, by the methods of
\hyperlink{bib:E}{E.}, p.~115 (before Theorem~11*), one could give a
simultaneous spectral representation for \(B,W,W^*\), and hence for
\(A,A^*\).  We prefer to use the definition in \hyperlink{bib:A}{A.},
p.~116, Definition~6, and argue as follows.  The operators \(A,A^*\) have
the same domains as \(B,C\), respectively, and hence the same domain as
one another; thus, in the sense of \hyperlink{bib:A}{A.}, p.~405,
Definition~5, they form a conjugate pair.  In the notation used there,
if \(E(\lambda)\) also denotes the resolution of the identity belonging
to \(B\), then
\[
  (A)''=(WB)''
  \subseteq (W,B)''
  =\bigl(W,W^*,\text{all }E(\lambda)\bigr)''.
\]
This set is Abelian, since it is the set generated by
\(W,W^*,E(\lambda)\); compare \hyperlink{bib:A}{A.}, p.~389, above.
Indeed, \(W,W^*\) commute with \(B\), and therefore with all
\(E(\lambda)\), by \hyperlink{bib:A}{A.}, p.~408, Theorem~13.

Conversely, suppose that \(A\) is normal.  Then the closed operators
\(A,A^*\) have the same domain (\hyperlink{bib:A}{A.}, p.~416,
Theorem~21).  If \(A^*Af\) and \(AA^*g\) are defined, then
\[
  (f,AA^*g)
  =(A^*f,A^*g)\footnotemark
  =(Af,Ag)
  =(A^*Af,g).
\]
\footnotetext{\label{fn:21}According to \hyperlink{bib:A}{A.}, at the
place just cited, whenever \(Ah\) is defined,
\(\lVert Ah\rVert=\lVert A^*h\rVert\), and hence
\((Ah,Ah)=(A^*h,A^*h)\).  From this one concludes, whenever \(Af\) and
\(Ag\) are defined, by considering
\(\frac{f\pm g}{2}\) and \(\frac{f\pm ig}{2}\), that
\((Af,Ag)=(A^*f,A^*g)\) as well.}
By the hypermaximality of \(AA^*\), it follows that \(AA^*f\) is defined
and equal to \(A^*Af\).  Thus \(AA^*\) is an extension of \(A^*A\); but
the latter too is hypermaximal, and consequently
\(A^*A=AA^*\).  Thus:

\par\smallskip
\phantomsection
\label{thm:8}
\noindent\textbf{Theorem 8.}
\textit{Let \(A\) be as in \hyperref[thm:2]{Theorems~2} and
\hyperref[thm:3]{3}, and let}
\[
  B=\sqrt{A^*A},
  \qquad
  C=\sqrt{AA^*}.
\]
\textit{For \(A\) to be normal it is necessary and sufficient that}
\[
  A^*A=AA^*,
\]
\textit{or, equivalently, that \(B=C\).  Since \(A^*A\) and \(AA^*\) are
both hypermaximal, it even suffices that one be an extension of the
other.}

% --- Source page 309 (PDF page 16) ---

Suppose that \(A,A^*\) have the same domain and that
\(\lVert Af\rVert=\lVert A^*f\rVert\) throughout it.  The same then holds
for \(B,C\).  The correspondence
\[
  Bf\Longleftrightarrow Cf
\]
maps \(\mathfrak B\) one-to-one, linearly, and isometrically onto
\(\mathfrak B'\).  It may therefore be extended in the same way to
\[
  \overline{\mathfrak B}=\mathfrak H-\mathfrak N,
  \qquad
  \overline{\mathfrak B}'=\mathfrak H-\mathfrak N';
\]
but \(\mathfrak N=\mathfrak N'\), and consequently
\(\overline{\mathfrak B}=\overline{\mathfrak B}'\).  Denote this operator
by \(V_0\).  It therefore maps
\(\overline{\mathfrak B}=\overline{\mathfrak B}'
=\mathfrak H-\mathfrak N\) onto itself.  Extend it to a linear operator
\(W_0\) which vanishes on \(\mathfrak N\).  As in
\hyperref[thm:6]{Theorem~6},
\[
  C=V_0B=W_0B,
  \qquad
  W_0^*W_0=W_0W_0^*=E.
\]

Taking adjoints gives \(C=BW_0^*\).  Thus \(W_0^*\), within
\(\mathfrak H-\mathfrak N\), carries the domain of \(C\) into that of
\(B\), that is, carries the domain of \(B\) into itself.  The same is true
of its inverse \(W_0\) on \(\mathfrak H-\mathfrak N\).  Hence \(Bf\) and
\(Cf=W_0Bf\), both of which lie in \(\mathfrak H-\mathfrak N\), belong
simultaneously to the common domain of \(B,C\).  In other words, \(B^2\)
and \(C^2\) have the same domain.

On that domain one always has
\[
  (B^2f,f)=(Bf,Bf)=\lVert Bf\rVert^2,
  \qquad
  (C^2f,f)=(Cf,Cf)=\lVert Cf\rVert^2.
\]
Since \((B^2f,f)=(C^2f,f)\), it follows that \(B^2=C^2\), and therefore
\(B=C\).  Thus:

\par\smallskip
\phantomsection
\label{thm:9}
\noindent\textbf{Theorem 9.}
\textit{The following is also a necessary and sufficient condition for
the normality of \(A\): \(A\) and \(A^*\) have the same domain, and
throughout it}
\[
  \lVert Af\rVert=\lVert A^*f\rVert.\footnotemark
\]
\footnotetext{\label{fn:22}In \hyperlink{bib:A}{A.}, p.~423, above, it
was shown that it is not sufficient for this equality to hold merely on
some everywhere dense set.}

\phantomsection
\label{subsec:II.2}
\noindent\textbf{2.}
We now turn to the notions of metric equality and metric extension of
operators defined in \hyperref[intro:4]{Introduction~4}.

From \hyperref[thm:6]{Theorem~6}, together with the argument used in the
proof of \hyperref[thm:9]{Theorem~9}, which shows that any two Hermitian,
definite, hypermaximal operators \(B,C\) that are metrically equal must
actually be equal (for the relation of \(B,C\) to \(A,A^*\) was not used
there), we immediately obtain:

\par\smallskip
\phantomsection
\label{thm:10}
\noindent\textbf{Theorem 10.}
\textit{Every closed linear \(*\)-operator \(A\) is metrically equal to
one and only one Hermitian, definite, hypermaximal operator \(B\), namely
the operator \(B\) of \hyperref[thm:6]{Theorem~6}:}
\[
  B=\sqrt{A^*A}.
\]

Every discontinuous \(A\) has proper metric extensions.  By
\hyperref[thm:10]{Theorem~10}, we may assume \(A\) to be Hermitian,
definite, and hypermaximal.  By \hyperlink{bib:E}{E.}, p.~108,
Theorem~49, \(A\) is a proper extension of a closed linear Hermitian
operator \(A'\), which can in fact be chosen in continuum-many distinct
ways.  Consequently \(A'^{*}\), which is a \(*\)-operator because \(A'\)
is a \(**\)-operator, is a proper extension of \(A\).  Thus:

% --- Source page 310 (PDF page 17) ---

\par\smallskip
\phantomsection
\label{thm:11}
\noindent\textbf{Theorem 11.}
\textit{Let \(A\) be as in \hyperref[thm:10]{Theorem~10}.  If \(A\) is
discontinuous, then it has proper metric extensions (indeed,
continuum-many distinct ones).}

If \(A\) is continuous, then its domain is closed; since it must be
everywhere dense, the domain is then equal to \(\mathfrak H\).  From this
and \hyperref[thm:11]{Theorem~11} it follows that:

\par\smallskip
\phantomsection
\label{thm:12}
\noindent\textbf{Theorem 12.}
\textit{Let \(A\) be as in \hyperref[thm:10]{Theorem~10}.  Then \(A\) is
continuous if and only if it is defined everywhere.}

As noted in \hyperref[intro:2]{Introduction~2}, this generalizes
Toeplitz's theorem%
\footnote{\label{fn:23}%
\hypertarget{bib:HellingerToeplitz1910}{%
\href{https://eudml.org/doc/158468}{E.~Hellinger and O.~Toeplitz,
``Grundlagen f\"ur eine Theorie der unendlichen Matrizen,''
\emph{Mathematische Annalen}, \textbf{69} (1910), pp.~289--330,
especially p.~321}}; compare also \hyperlink{bib:E}{E.}, pp.~107--108,
Theorem~48, \(\gamma\).}%
\translatornote{Source note~23 prints the year as 1911; the cited volume
and article were published in 1910.}.  This theorem can incidentally also
be proved directly.

There are also very general theorems on the simultaneous metric extension
of several operators, for example of operators whose domains have no
points in common (compare \hyperlink{bib:M}{M.}, p.~232, Theorem~18).
The author intends to discuss these shortly.

% --- End inlined translation_pages_15_17_backup.tex ---

\end{document}
